% ===== PRÉAMBULE CANONIQUE — à coller VERBATIM en tête de CHAQUE .tex =====
\documentclass[11pt,a4paper]{article}
\usepackage[utf8]{inputenc}\usepackage[T1]{fontenc}\usepackage{lmodern}
\usepackage[french]{babel}
\usepackage{amsmath,amssymb,mathtools}\usepackage{icomma}
\usepackage[version=4]{mhchem}          % \ce{...} : équations & formules
\usepackage{siunitx}\sisetup{locale=FR} % \qty{}{}, \si{}, \ang{}
\usepackage{chemfig}                    % \chemfig{...} : structures développées
\usepackage{xcolor}\usepackage{colortbl}
\usepackage{tikz}\usepackage{pgfplots}\pgfplotsset{compat=1.18}
\usepackage[most]{tcolorbox}\usepackage{enumitem}\usepackage{array,booktabs}
\usepackage[a4paper,margin=2.2cm]{geometry}
\usetikzlibrary{arrows.meta,calc,angles,quotes,positioning,patterns,backgrounds,shapes.geometric}
\definecolor{primary}{RGB}{222,49,99}\definecolor{primarydark}{RGB}{150,22,55}
\definecolor{defcol}{RGB}{0,114,178}\definecolor{methcol}{RGB}{52,92,148}
\definecolor{excol}{RGB}{0,135,90}\definecolor{attncol}{RGB}{200,40,55}
\tcbset{boxbase/.style={enhanced,breakable,boxrule=0.9pt,arc=2.5pt,
  left=3mm,right=3mm,top=2.3mm,bottom=2.3mm,fonttitle=\bfseries,coltitle=white}}
% --- boîtes TUTORIEL (noms EXACTS, ne pas renommer) ---
\newtcolorbox{cle}[1][]{boxbase,colback=defcol!6,colframe=defcol,
  title={\textsc{À retenir}\ifx\relax#1\relax\relax\else\textmd{ : \itshape #1}\fi}}
\newtcolorbox{methode}[1][]{boxbase,colback=methcol!7,colframe=methcol,sharp corners,
  title={\textsc{Méthode}\ifx\relax#1\relax\relax\else\textmd{ --- \itshape #1}\fi}}
\newtcolorbox{exemplebox}[1][]{boxbase,colback=excol!5,colframe=excol,
  title={\textsc{Exemple}\ifx\relax#1\relax\relax\else\textmd{ : \itshape #1}\fi}}
\newtcolorbox{piege}[1][]{boxbase,colback=attncol!6,colframe=attncol,
  title={\textsc{Piège}\ifx\relax#1\relax\relax\else\textmd{ : \itshape #1}\fi}}
\newtcolorbox{remarque}[1][]{boxbase,colback=black!4,colframe=black!45,
  title={\textsc{Remarque}\ifx\relax#1\relax\relax\else\textmd{ : \itshape #1}\fi}}
% --- boîtes EXERCICE / CORRIGÉ (noms EXACTS, auto-comptées) ---
\newtcolorbox[auto counter]{exo}[1][]{boxbase,colback=primary!4,colframe=primary,
  title={\textsc{Exercice}~\thetcbcounter\ifx\relax#1\relax\relax\else\textmd{ --- \itshape #1}\fi}}
\newtcolorbox[auto counter]{corrige}[1][]{boxbase,colback=excol!4,colframe=excol,
  title={\textsc{Corrigé}~\thetcbcounter\ifx\relax#1\relax\relax\else\textmd{ --- \itshape #1}\fi}}
\newcommand{\R}{\mathbb{R}}\newcommand{\N}{\mathbb{N}}\newcommand{\Z}{\mathbb{Z}}\newcommand{\ds}{\displaystyle}
% (un \usepackage{fancyhdr}+en-tête est possible ; il est ignoré côté web)
% =========================================================================
\begin{document}
\sisetup{per-mode=symbol}
\begin{exo}[juste ou fidèle ?]
On pèse quatre fois un objet de masse \emph{vraie} $\qty{10,00}{\gram}$. Pour chaque série de résultats,
dis si la mesure est \textbf{juste}, \textbf{fidèle}, les deux, ou ni l'un ni l'autre.
\begin{enumerate}[label=\alph*)]
  \item $10,01\,;\ 9,99\,;\ 10,00\,;\ 10,02$
  \item $10,51\,;\ 10,49\,;\ 10,50\,;\ 10,52$
  \item $9,5\,;\ 10,4\,;\ 9,8\,;\ 10,3$
  \item $11,2\,;\ 10,8\,;\ 12,0\,;\ 9,3$
\end{enumerate}
\end{exo}

\begin{exo}[incertitude absolue sous-entendue]
Donne l'incertitude absolue $\Delta x$ suggérée par le dernier chiffre significatif de chaque mesure.
\begin{enumerate}[label=\alph*)]
  \item $\qty{13,02}{\gram}$
  \item $\qty{4,5}{\milli\litre}$
  \item $\qty{250,0}{\gram}$
  \item $\qty{1,235}{\second}$
\end{enumerate}
\end{exo}

\begin{exo}[encadrer la vraie valeur]
À partir de l'écriture $x\pm\Delta x$, donne l'intervalle dans lequel se trouve la vraie valeur.
\begin{enumerate}[label=\alph*)]
  \item $(13,02\pm0,01)\ \si{\gram}$
  \item $(5,60\pm0,05)\ \si{\centi\metre}$
  \item $(20,96\pm0,01)\ \si{\milli\litre}$
\end{enumerate}
\end{exo}

\begin{exo}[incertitude relative]
Calcule l'incertitude relative $\dfrac{\Delta x}{x}$, exprimée en pourcentage (à $2$ CS).
\begin{enumerate}[label=\alph*)]
  \item $(13,02\pm0,01)\ \si{\gram}$
  \item $(4,5\pm0,1)\ \si{\centi\metre}$
  \item $(250,0\pm0,1)\ \si{\gram}$
\end{enumerate}
\end{exo}

\end{document}
